\documentclass[11pt,a4paper]{article}

\usepackage[T1]{fontenc}
\usepackage[utf8]{inputenc}
\usepackage[provide=*,french]{babel}
\usepackage[a4paper,margin=2.3cm]{geometry}
\usepackage{amsmath,amssymb}
\usepackage{booktabs,longtable,array}
\usepackage{hyperref}
\usepackage{enumitem}
\usepackage{listings}
\usepackage{xcolor}

\hypersetup{
  colorlinks=true,
  linkcolor=blue,
  urlcolor=blue,
  pdftitle={Conception technique - Inventaires completude representativite (ICR)}
}

\lstset{
  basicstyle=\ttfamily\small,
  breaklines=true,
  frame=single,
  backgroundcolor=\color{gray!6}
}

\title{Conception Technique\\
\large Inventaires\_completude\_representativite}
\author{Boite Eddy}
\date{\today}

\hbadness=10000
\hfuzz=700pt
\begin{document}
\sloppy
\maketitle

\begin{abstract}
Ce document présente la conception technique du pipeline d'analyse de complétude et de
représentativité d'inventaires fongiques \texttt{Inventaires\_completude\_representativite.R}
(v1.4). La description adopte un cadre d'ingénierie reproductible : architecture, contrats de
données, configuration runtime, robustesse opérationnelle et livrables produits. L'objectif est
de fournir une référence techniquement auditable et réutilisable pour l'analyse d'inventaires
fongiques par site.
\end{abstract}
\tableofcontents
\newpage

\section{Structuration IMRaD adaptée}
Dans une perspective de publication scientifique, la lecture du document peut
être organisée selon le schéma IMRaD, adapté à une documentation technique :
\begin{itemize}[leftmargin=1.2cm]
  \item \textbf{Introduction} : objectifs techniques et périmètre du script.
  \item \textbf{Méthodes} : architecture, contrats d'entrée/sortie, configuration runtime, pipeline.
  \item \textbf{Résultats} : organisation des sorties, livrables et journalisation.
  \item \textbf{Discussion} : limites connues et robustesse.
\end{itemize}

\section{Positionnement méthodologique}
Le pipeline relève d'une approche \textit{data-centric} appliquée à l'écologie descriptive de
terrain : les indicateurs de complétude et de représentativité sont construits par accumulation
successive des observations au fil des visites, puis consolidés par des modèles d'ajustement et
des métriques de pertinence agrégées. Deux objectifs sont poursuivis simultanément :
\begin{itemize}[leftmargin=1.2cm]
  \item \textbf{validité analytique} (qualité des ajustements de courbes d'accumulation, robustesse des indicateurs TEE/Ir),
  \item \textbf{validité opérationnelle} (stabilité des sorties, traçabilité et exploitabilité décisionnelle par site).
\end{itemize}

\section{Cadre théorique des choix statistiques}
Le dispositif s'appuie sur la théorie des courbes d'accumulation d'espèces
(\textit{species accumulation curves}) et sur une estimation empirique du taux d'espèces
rares. Cette combinaison répond à une double exigence : décrire fidèlement la dynamique de
découverte des espèces au fil des visites, et estimer une richesse asymptotique théorique
permettant de quantifier objectivement l'effort d'échantillonnage restant.

Dans ce cadre, la sélection des métriques privilégie :
\begin{itemize}[leftmargin=1.2cm]
  \item la \textbf{comparabilité inter-sites} (Ir et complétude bornés dans $[0,1]$),
  \item la \textbf{sensibilité à la stabilisation de la courbe} (pente terminale, part de nouveauté en fin d'inventaire),
  \item la \textbf{robustesse à l'absence de données optionnelles} (dégradation contrôlée si \texttt{placette}, \texttt{minpack.lm} ou \texttt{vegan} sont indisponibles),
  \item la \textbf{traçabilité des décisions} (logs horodatés et manifeste final des livrables).
\end{itemize}

\section{Objectif technique}
Ce document décrit l'architecture technique du script
\texttt{scripts/Inventaires\_completude\_representativite.R} (v1.4).
Il précise les contrats d'entrée/sortie, la configuration runtime
et les limites opérationnelles.

\section{Architecture logicielle}
\subsection{Blocs principaux}
\begin{enumerate}
  \item Initialisation et validation de la configuration (\texttt{validate\_config()})
  \item Résolution du fichier d'entrée (\texttt{resolve\_input\_file()})
  \item Lecture avec détection automatique du délimiteur (\texttt{read\_delim\_auto()})
  \item Audit de conformité CSV et export du rapport (\texttt{audit\_csv\_conformity()}, \texttt{export\_csv\_conformity\_report()})
  \item Préparation et nettoyage des données (\texttt{prepare\_data()})
  \item Analyse par site (\texttt{analyze\_site()} appliquée par \texttt{purrr::map} sur les groupes de site)
  \item Consolidation multi-sites et graphiques comparatifs
  \item Manifeste final des livrables (\texttt{log\_output\_manifest()})
\end{enumerate}

\subsection{Fonctions clés}
\begin{itemize}
  \item \textbf{Configuration et logging} : \texttt{get\_embedded\_config()}, \texttt{setup\_logging()}, \texttt{log\_message()}, \texttt{log\_info()}, \texttt{log\_debug()}, \texttt{log\_warning()}, \texttt{log\_header()}, \texttt{log\_section()}
  \item \textbf{R\'esolution et validation} : \texttt{resolve\_path()}, \texttt{resolve\_input\_file()}, \texttt{validate\_config()}, \texttt{check\_required\_packages()}
  \item \textbf{Lecture et audit} : \texttt{read\_delim\_auto()}, \texttt{audit\_csv\_conformity()}, \texttt{export\_csv\_conformity\_report()}
  \item \textbf{Pr\'eparation} : \texttt{prepare\_data()}
  \item \textbf{Analyses univari\'ees} : \texttt{calc\_cumulative\_metrics()}, \texttt{calc\_tee\_ir()}, \texttt{calc\_species\_frequency()}, \texttt{calc\_spatial\_occupancy()}
  \item \textbf{Mod\'elisation} : \texttt{fit\_linear\_model()}, \texttt{fit\_hyperbolic\_model()} (asymptote extraite via \texttt{extract\_model\_stats()})
  \item \textbf{M\'etriques avanc\'ees} : \texttt{calc\_scientific\_metrics()}, \texttt{run\_ca()}
  \item \textbf{Orchestration} : \texttt{analyze\_site()}, \texttt{run\_analysis()}, \texttt{log\_output\_manifest()}, \texttt{ensure\_thematic\_dirs()}
  \item \textbf{Visualisation} : \texttt{theme\_project()}, \texttt{PALETTE\_PRIMARY} (5 teintes CVD-friendly), \texttt{save\_plot()}, fonctions \texttt{plot\_*()}
\end{itemize}

\subsection{Points d'entrée}
\begin{itemize}[leftmargin=1.2cm]
  \item Le lancement automatique est conditionné par le flag \texttt{AUTO\_RUN} (piloté par \texttt{INVENTAIRES\_AUTO\_RUN}) : en session non interactive (\texttt{Rscript}), \texttt{run\_analysis(CONFIG)} est appelé directement ; en session interactive, un message informe que la fonction est prête à être invoquée manuellement.
  \item Contrairement au script CHEGD de référence, aucun bloc \texttt{tryCatch} global n'encapsule \texttt{run\_analysis()} : une erreur non capturée localement (ex. \texttt{stop()} de \texttt{validate\_config()} ou \texttt{prepare\_data()}) interrompt l'exécution complète (voir \S\,\ref{sec:robustesse}).
\end{itemize}

\section{Dépendances techniques}
\subsection{Packages R}
\begin{itemize}[leftmargin=1.2cm]
  \item Requis : \texttt{dplyr}, \texttt{tidyr}, \texttt{ggplot2}, \texttt{purrr}, \texttt{readr}, \texttt{stringr}, \texttt{forcats}, \texttt{tibble}, \texttt{scales}, \texttt{gridExtra}.
  \item Optionnels : \texttt{minpack.lm} (\texttt{nlsLM}, modèle hyperbolique), \texttt{vegan} (\texttt{cca}, CA/AFC).
\end{itemize}

\subsection{Comportement d'installation}
\texttt{check\_required\_packages()} vérifie la disponibilité de chaque package requis via
\texttt{requireNamespace()}. En cas d'absence, le script s'arrête avec une erreur explicite et
une commande \texttt{install.packages()} suggérée (pas d'installation automatique). Les
packages optionnels sont détectés via les indicateurs \texttt{HAS\_VEGAN}/\texttt{HAS\_MINPACK}
sans bloquer l'exécution : les fonctionnalités correspondantes sont simplement désactivées
avec un message \texttt{info} explicite.

\section{Contrats d'entrée et de sortie}

\subsection{Contrat d'entrée}
\texttt{data/observations.csv} (ou \texttt{.txt}/\texttt{.tsv}) avec colonnes obligatoires \texttt{site,date,visite\_id,espece} et optionnelle \texttt{placette}.

Règles : visite distincte \texttt{date + visite\_id} (et \texttt{site} en multi-sites), parsing date strict, déduplication de la clé observation.

\subsection{Configuration runtime (CONFIG)}
\begin{longtable}{|p{4.2cm}|p{3.1cm}|p{7.0cm}|}
\hline
\textbf{Cl\'e} & \textbf{D\'efaut} & \textbf{Description} \\
\hline
\texttt{input\_file} & \path{data/observations.csv} & Fichier source \\
\hline
\texttt{output\_dir} & \texttt{results/ICR} & R\'epertoire de sortie \\
\hline
\texttt{date\_format} & \texttt{\%Y-\%m-\%d} & Parsing des dates \\
\hline
\texttt{csv\_strict\_mode} & \texttt{TRUE} & Arr\^et si CSV non conforme \\
\hline
\texttt{csv\_allow\_extra\_cols} & \texttt{FALSE} & Tol\'erance des colonnes suppl\'ementaires \\
\hline
\texttt{csv\_required\_cols} & \path{site,date,visite_id,espece} & Sch\'ema obligatoire \\
\hline
\texttt{csv\_optional\_cols} & \texttt{placette} & Sch\'ema optionnel \\
\hline
\texttt{freq\_breaks} & \texttt{0, 0.10, 0.25,} & Séquence des intervalles de fréquence \\
 & \texttt{0.50, 0.75, 1.00} & (proportions de visites, dans $[0,1]$) \\
\hline
\texttt{freq\_labels} & \texttt{exceptionnelle,} & Étiquettes des intervalles \\
 & \texttt{très\_rare, occasionnelle,} & \\
 & \texttt{fréquente, constante} & \\
\hline
\texttt{make\_ca} & \texttt{TRUE} & Active la CA/AFC si possible \\
\hline
\texttt{min\_visits\_for\_model} & \texttt{5} & Seuil de mod\'elisation \\
\hline
\texttt{width} & \texttt{10} & Largeur graphiques PNG (pouces) \\
\hline
\texttt{height} & \texttt{6} & Hauteur graphiques PNG (pouces) \\
\hline
\texttt{dpi} & \texttt{300} & R\'esolution graphiques PNG (points/pouce) \\
\hline
\end{longtable}

\noindent
Seules cinq clés sont effectivement surchargeables via variables d'environnement dans
\texttt{get\_embedded\_config()} : \texttt{input\_file}, \texttt{output\_dir}, \texttt{date\_format},
\texttt{csv\_strict\_mode} et \texttt{csv\_allow\_extra\_cols}. Les clés \texttt{freq\_breaks},
\texttt{freq\_labels}, \texttt{make\_ca}, \texttt{min\_visits\_for\_model}, \texttt{width},
\texttt{height} et \texttt{dpi} sont actuellement des constantes du code source (littéraux R) :
leur modification requiert d'éditer la fonction \texttt{get\_embedded\_config()} elle-même,
malgré le principe de configuration centralisée. La configuration effective est affichée
dans l'en-tête du script via la fonction \texttt{print\_startup\_header()}.

Variables d'environnement effectivement lues : \texttt{INVENTAIRES\_INPUT\_FILE},
\texttt{INVENTAIRES\_OUTPUT\_DIR}, \texttt{INVENTAIRES\_DATE\_FORMAT},
\texttt{INVENTAIRES\_CSV\_STRICT}, \texttt{INVENTAIRES\_CSV\_ALLOW\_EXTRA\_COLS} et
\texttt{INVENTAIRES\_AUTO\_RUN}. Les autres options reprennent les cl\'es de configuration.

\subsection{Modes d'exécution et flags internes}
Le script expose des flags globaux :
\begin{itemize}
  \item \texttt{DEBUG\_MODE} (d\'efaut \texttt{FALSE}) : warnings affich\'es en continu et logs d\'etaill\'es.
  \item \texttt{BENCHMARK\_MODE} (d\'efaut \texttt{FALSE}) : chronom\'etrage des blocs de calcul.
  \item \texttt{CLEAN\_ENVIRONMENT} (d\'efaut \texttt{FALSE}) : purge de l'environnement avant lancement (mode expert).
\end{itemize}

Le lancement automatique est contrôlé via \texttt{INVENTAIRES\_AUTO\_RUN} (défaut \texttt{TRUE}).

\subsection{Audit CSV}
Sorties d'audit (sous \texttt{results/ICR/00\_data\_quality/}) :
\begin{itemize}
  \item \texttt{ICR\_00\_csv\_conformite\_report.csv}
  \item \texttt{ICR\_00\_csv\_conformite\_problems.csv} (si anomalies de parsing)
\end{itemize}

En mode strict (\texttt{csv\_strict\_mode = TRUE}, défaut), toute non-conformité (colonne
manquante, colonne supplémentaire non tolérée, problème de parsing) est bloquante
(\texttt{stop()}). En mode tolérant (\texttt{FALSE}), un \texttt{log\_warning()} est émis et le
pipeline se poursuit sur les données telles que lues (voir \S\,\ref{sec:pipeline}).

\section{Modèle de données technique}

\subsection{Schéma logique des observations}
\begin{longtable}{|p{2.4cm}|p{2.6cm}|p{9.3cm}|}
\hline
\textbf{Champ} & \textbf{Type} & \textbf{Rôle} \\
\hline
\texttt{site} & caractère & Clé de groupement principale ; détermine le découpage \texttt{analyze\_site()} et le sous-répertoire \texttt{sites/<slug>/}. \\
\hline
\texttt{date} & date (parsée selon \texttt{date\_format}) & Ordonnancement chronologique des visites ; composant de la clé de déduplication. \\
\hline
\texttt{visite\_id} & caractère & Combiné à \texttt{date} pour définir une visite distincte (\texttt{count\_distinct\_visits()}). \\
\hline
\texttt{espece} & caractère (normalisé via \texttt{stringr::str\_squish()}) & Unité d'accumulation pour toutes les courbes et indices (TEE, Ir, fréquence). \\
\hline
\texttt{placette} & caractère, optionnel & Active \texttt{calc\_spatial\_occupancy()} et \texttt{run\_ca()} si renseignée (non totalement \texttt{NA}/vide) ; sinon les sorties correspondantes sont omises (statut \texttt{OPTIONNEL\_NON\_GENERE}). \\
\hline
\end{longtable}

Clé de déduplication appliquée par \texttt{prepare\_data()} : \texttt{(site, placette, date, visite\_id, espece)}.

\subsection{Artefacts intermédiaires}
\begin{itemize}[leftmargin=1.2cm]
  \item \textbf{Données préparées} : table nettoyée/dédupliquée, exportée telle quelle (\texttt{ICR\_donnees\_preparees.csv}).
  \item \textbf{Table des visites} (\texttt{visits\_tbl}) : visites distinctes \texttt{date + visite\_id}, triées chronologiquement, indexées par \texttt{visite\_index}.
  \item \textbf{Table cumulée} (\texttt{df\_cum}) : une ligne par visite avec \texttt{nb\_especes\_trouvees}, \texttt{nb\_nouvelles}, \texttt{nb\_cumule}.
  \item \textbf{Table TEE/Ir} (\texttt{df\_tee}) : une ligne par visite avec \texttt{tee} et \texttt{ir}.
  \item \textbf{Table de fréquences} (\texttt{freq\_tbl}) : une ligne par espèce avec \texttt{prop\_visites} et classe qualitative.
  \item \textbf{Table d'occupation spatiale} (\texttt{occ\_tbl}, optionnelle) : une ligne par espèce avec \texttt{prop\_placettes}.
  \item \textbf{Statistiques de modèles} (\texttt{model\_stats}) : R², équation et asymptote pour les modèles linéaire et hyperbolique.
  \item \textbf{Métriques scientifiques} (\texttt{scientific\_metrics}) : agrégats par site issus de \texttt{calc\_scientific\_metrics()}.
  \item \textbf{Consolidations globales} : \texttt{site\_summaries} et \texttt{site\_metrics}, obtenues par \texttt{purrr::map()} + \texttt{dplyr::bind\_rows()} sur les résultats de chaque site.
\end{itemize}

\section{Pipeline technique détaillé}
\label{sec:pipeline}

\subsection{Phase A --- Ingestion et normalisation}
\begin{itemize}[leftmargin=1.2cm]
  \item \texttt{resolve\_input\_file()} : tente le chemin configuré (\texttt{input\_file}/\texttt{INVENTAIRES\_INPUT\_FILE}) ; à défaut, recherche un fichier \texttt{observations.\{csv,txt,tsv\}} sous \texttt{data/} ; \texttt{stop()} explicite si aucun candidat trouvé.
  \item \texttt{read\_delim\_auto()} : détecte le délimiteur en comptant les occurrences de tabulation, point-virgule et virgule sur la première ligne, puis lit via \texttt{readr::read\_delim()} (locale UTF-8, toutes colonnes en caractère par défaut via \texttt{build\_csv\_colspec()}) ; les problèmes de parsing sont capturés via \texttt{readr::problems()}.
  \item \texttt{audit\_csv\_conformity()} : calcule colonnes manquantes/supplémentaires, nombre de problèmes de parsing et statut \texttt{CONFORME}/\texttt{NON\_CONFORME} ; \texttt{export\_csv\_conformity\_report()} écrit le rapport dans \texttt{00\_data\_quality/}.
  \item \texttt{prepare\_data()} : vérifie la présence des colonnes obligatoires, normalise les champs texte, parse les dates, \texttt{stop()} explicite si valeurs manquantes/vides sur \texttt{site}/\texttt{visite\_id}/\texttt{espece}/\texttt{date}, ajoute \texttt{placette = NA\_character\_} si absente, déduplique sur la clé composée et journalise les effectifs avant/après.
\end{itemize}

\subsubsection{Critères de qualité de données}
\begin{itemize}[leftmargin=1.2cm]
  \item Complétude des colonnes obligatoires (arrêt si absentes).
  \item Validité typologique des dates (arrêt si non parsées selon \texttt{date\_format}).
  \item Traçabilité des doublons supprimés (nombre de lignes avant/après déduplication journalisé).
\end{itemize}

\subsection{Phase B --- Indicateurs cumulés et de représentativité}

\subsubsection{Courbe temps-espèces}
\texttt{calc\_cumulative\_metrics()} trie les visites chronologiquement (\texttt{date}, \texttt{visite\_id}), leur attribue un \texttt{visite\_index}, puis parcourt chaque visite $i$ pour calculer :
\begin{itemize}[leftmargin=1.2cm]
  \item \texttt{nb\_especes\_trouvees} : nombre d'espèces observées lors de la visite $i$,
  \item \texttt{nb\_nouvelles} : cardinal de la différence ensembliste entre les espèces de la visite $i$ et l'ensemble des espèces déjà vues (\texttt{setdiff()}),
  \item \texttt{nb\_cumule} : cardinal de l'union cumulative des espèces vues jusqu'à $i$ (\texttt{union()}).
\end{itemize}

\subsubsection{Taux d'espèces exclusives (TEE) et indice de représentativité (Ir)}
\texttt{calc\_tee\_ir()} considère, pour chaque visite $i$, l'ensemble des espèces vues sur les
visites $1$ à $i$ :
$$ TEE_i = \frac{s_{1,i}}{s_{total,i}} \qquad Ir_i = 1 - TEE_i $$
où $s_{1,i}$ est le nombre d'espèces vues exactement une fois parmi les visites $1..i$, et
$s_{total,i}$ le nombre d'espèces distinctes vues parmi les visites $1..i$. Repères de
lecture utilisés dans les graphiques : $Ir < 0.60$ (faible), $0.60 \le Ir < 0.80$ (moyenne),
$Ir \ge 0.80$ (bonne).

\subsection{Phase C --- Modélisation et complétude}
\begin{itemize}[leftmargin=1.2cm]
  \item \textbf{Modèle linéaire} : \texttt{fit\_linear\_model()} ajuste \texttt{stats::lm(nb\_cumule \textasciitilde{} visite\_index)}.
  \item \textbf{Modèle hyperbolique} (si \texttt{HAS\_MINPACK}) : \texttt{fit\_hyperbolic\_model()} ajuste, via \texttt{minpack.lm::nlsLM} (Levenberg-Marquardt, 500 itérations max), l'équation
  $$ S(t) = a - \frac{1}{bt + c} $$
  avec valeurs de départ $a_0 = 1.20 \times \max(\texttt{nb\_cumule})$, $b_0 = 10^{-4}$, $c_0 = 10^{-3}$.
  \item \textbf{Extraction des statistiques} : \texttt{extract\_model\_stats()} calcule $R^2 = 1 - SS_{res}/SS_{tot}$ ; pour le modèle hyperbolique, l'asymptote verticale est $a$ et l'asymptote horizontale (abscisse théorique d'origine) est $-c/b$.
  \item \textbf{Complétude} : $\texttt{completude} = \texttt{nb\_cumule}_{max} / a$ (asymptote hyperbolique), non calculée (\texttt{NA}) si $a$ est absent ou $\le 0$. Repères : $< 0.70$ (insuffisante), $0.70$--$0.90$ (intermédiaire), $\ge 0.90$ (avancée).
  \item \textbf{Condition d'activation} : modélisation déclenchée uniquement si \texttt{nrow(df\_cum) >= config\$min\_visits\_for\_model} (défaut $5$).
\end{itemize}

\subsection{Phase D --- Fréquences, occupation spatiale et CA/AFC}
\begin{itemize}[leftmargin=1.2cm]
  \item \texttt{calc\_species\_frequency()} : $\texttt{prop\_visites} = \texttt{nb\_visites} / \texttt{n\_visites\_totales}$, discrétisé via \texttt{cut(freq\_breaks, freq\_labels, include.lowest = TRUE, right = TRUE)}.
  \item \texttt{calc\_spatial\_occupancy()} : requiert une colonne \texttt{placette} non entièrement \texttt{NA}/vide ; $\texttt{prop\_placettes} = \texttt{nb\_placettes} / \texttt{n\_placettes\_totales}$ ; retourne \texttt{NULL} sinon.
  \item \texttt{run\_ca()} : requiert \texttt{HAS\_VEGAN} et des données de placette ; construit une matrice binaire placette $\times$ espèce (\texttt{tidyr::pivot\_wider()}), exige au moins 3 colonnes d'espèces, puis appelle \texttt{vegan::cca()}.
\end{itemize}

\subsection{Phase E --- Métriques de pertinence scientifique}
\texttt{calc\_scientific\_metrics()} calcule, par site :
\begin{itemize}[leftmargin=1.2cm]
  \item \textbf{Pente terminale} (\texttt{slope\_terminal\_cumul}) : pente d'une régression linéaire sur la fenêtre terminale de $k = \min(5, n_{visites})$ points de la courbe cumulée.
  \item \textbf{Part de nouveauté terminale} (\texttt{tail\_novelty\_share}) : $\sum \texttt{nb\_nouvelles}_{fen\hat{e}tre} / \sum \texttt{nb\_nouvelles}_{total}$.
  \item \textbf{TEE/Ir finaux} : dernière valeur (\texttt{dplyr::last()}) des séries \texttt{tee}/\texttt{ir}.
  \item \textbf{$R^2$ retenu} : celui du modèle hyperbolique si disponible, sinon celui du modèle linéaire.
  \item \textbf{Corrélation temporel/spatial} (\texttt{spearman\_temporal\_spatial}) : corrélation de Spearman (\texttt{stats::cor(method = "spearman")}) entre \texttt{prop\_visites} et \texttt{prop\_placettes} pour les espèces communes à \texttt{freq\_tbl} et \texttt{occ\_tbl}, calculée uniquement si $n \ge 3$ espèces jointes.
  \item \textbf{Taux de discordance} (\texttt{discordance\_rate}) : proportion d'espèces pour lesquelles $|\texttt{prop\_visites} - \texttt{prop\_placettes}| > 0.50$.
  \item \textbf{Score de pertinence} (\texttt{score\_pertinence}) : moyenne pondérée
  $$ score = 0.35 \cdot Ir_{final} + 0.35 \cdot completude + 0.20 \cdot (1 - \texttt{tail\_novelty\_share}) + 0.10 \cdot R^2_{mod\grave{e}le} $$
  les composantes manquantes (\texttt{NA}) sont exclues et les poids restants renormalisés.
  \item \textbf{Classes qualitatives} : \texttt{class\_ir} ($<0.60$ faible, $<0.80$ moyenne, $\ge 0.80$ bonne) et \texttt{class\_completude} ($<0.70$ insuffisante, $<0.90$ intermédiaire, $\ge 0.90$ avancée).
\end{itemize}

\subsection{Phase F --- Visualisation}
Toutes les figures partagent \texttt{theme\_project()} et la palette \texttt{PALETTE\_PRIMARY}
(bleu \texttt{\#0072B2}, orange \texttt{\#E69F00}, vert \texttt{\#009E73}, rouge \texttt{\#D55E00},
violet \texttt{\#CC79A7} --- palette adaptée au daltonisme). Fonctions principales :
\texttt{plot\_richness\_over\_time()}, \texttt{plot\_time\_species\_curve()} (annotation de
l'asymptote hyperbolique), \texttt{plot\_tee\_ir()} (deux panneaux via
\texttt{gridExtra::grid.arrange()}, zones de couleur $<0.60$/$0.60$--$0.80$/$\ge 0.80$),
\texttt{plot\_frequency\_hist()}, \texttt{plot\_spatial\_vs\_temporal()} (conditionnelle),
\texttt{plot\_scientific\_metrics\_site/heatmap/score()}. \texttt{save\_plot()} exporte via
\texttt{ggplot2::ggsave()} (format PNG dans tous les appels du script, \texttt{width}/\texttt{height}/\texttt{dpi} configurés), avec \texttt{tryCatch} et \texttt{log\_warning()} en cas d'échec d'écriture.

\section{Résultats}

\subsection{Logging systématisé}
Chaque exécution génère un fichier log horodaté \textbf{sous \texttt{logs/} à la racine du
projet} (\texttt{PROJECT\_DIR/logs/}, où \texttt{PROJECT\_DIR} est le répertoire parent du
dossier \texttt{scripts/}) --- et non sous \texttt{results/ICR/} comme les autres livrables,
puisque \texttt{setup\_logging()} est appelé avec \texttt{base\_dir = PROJECT\_DIR}.

\subsubsection{Fonctions de logging}
\begin{itemize}
  \item \texttt{setup\_logging(base\_dir, prefix)} : Initialise l'environnement .log\_env et crée le répertoire logs/ ; appelée au démarrage
  \item \texttt{log\_message(level, ...)} : Fonction centrale avec gestion intelligente des format strings
  \item \texttt{log\_info(...)} : Wrapper niveau INFO
  \item \texttt{log\_debug(...)} : Wrapper niveau DEBUG
  \item \texttt{log\_warning(...)} : Wrapper niveau WARN (affiche aussi en stderr)
  \item \texttt{log\_header(script\_name, script\_version)} : Affiche l'en-t\^ete du script avec métadonnées
  \item \texttt{log\_section(title)} : Affiche un titre de section avec délimiteurs
\end{itemize}

\subsubsection{Format et contenu}
\begin{itemize}
  \item \textbf{Format des entrées} : \texttt{[YYYY-MM-DD HH:MM:SS] [LEVEL] message}
  \item \textbf{Nom du fichier} : \texttt{PROJECT\_DIR/logs/ICR\_YYYYMMDD\_HHMM.log}
  \item \textbf{Contenu} : En-tête du script, métadonnées (timestamp, R version, working dir), progression du pipeline, manifeste final des livrables
\end{itemize}

\subsection{Organisation thématique des sorties}
Les résultats sont organisés selon une hiérarchie prédéfinie sous \texttt{results/ICR/}
(créée par \texttt{ensure\_thematic\_dirs()}) ; les noms exacts des répertoires réservés
(\texttt{04\_modeles}, \texttt{05\_spatial}) sont ceux définis dans le code source :

\begin{verbatim}
results/ICR/
  +-- 00_data_quality/
  |   +-- ICR_00_csv_conformite_report.csv
  |   +-- ICR_00_csv_conformite_problems.csv (si anomalies)
  |   `-- ICR_donnees_preparees.csv
  +-- 01_richesse/ (réservé pour extensions)
  +-- 02_completude_repr/ (réservé)
  +-- 03_frequences/ (réservé)
  +-- 04_modeles/ (réservé)
  +-- 05_spatial/ (réservé)
  +-- 06_metrics/
  |   +-- ICR_metrics_pertinence_tous_sites.csv
  |   +-- ICR_metrics_pertinence_heatmap_sites.png
  |   `-- ICR_metrics_pertinence_score_sites.png
  +-- comparisons/
  |   +-- ICR_resume_tous_sites.csv
  |   +-- ICR_comparaison_completude_sites.png
  |   `-- ICR_comparaison_ir_sites.png
  `-- sites/
      +-- <site_name>/
      |   +-- ICR_01_courbe_temps_especes.csv
      |   +-- ICR_02_tee_ir.csv
      |   +-- ICR_03_frequence_especes.csv
      |   +-- ICR_04_resume_site.csv
      |   +-- ICR_07_metrics_pertinence.csv
      |   +-- ICR_01_richesse_par_visite_et_cumul.png
      |   +-- ICR_02_courbe_temps_especes_hyperbole.png
      |   +-- ICR_03_tee_ir.png
      |   +-- ICR_04_histogramme_frequences.png
      |   +-- ICR_07_metrics_pertinence_dashboard.png
      |   +-- ICR_05_modeles.csv (si modèles activés)
      |   +-- ICR_06_occupation_spatiale.csv (si placette)
      |   +-- ICR_05_temporel_vs_spatial.png (conditionnelle)
      |   `-- ICR_06_CA_placettes_especes.png (si vegan + placette)
      `-- <other_site_name>/
          `-- ...

PROJECT_DIR/logs/
  `-- ICR_YYYYMMDD_HHMM.log
\end{verbatim}

\subsection{Journalisation et intégrité des sorties}
La journalisation est réalisée en console via \texttt{log\_info()}, \texttt{log\_debug()} et \texttt{log\_warning()}.

En fin de pipeline, \texttt{log\_output\_manifest()} contr\^ole les livrables avec statuts :
\texttt{OK}, \texttt{MANQUANT}, \texttt{OPTIONNEL\_NON\_GENERE}.

Le script supprime automatiquement tout \texttt{Rplots.pdf} parasite généré pendant l'exécution (niveau site et global).

\subsection{Livrables}
\subsubsection{Global}
\begin{itemize}
  \item \textbf{Qualité des données} (\texttt{00\_data\_quality/}) :
  \texttt{ICR\_00\_csv\_conformite\_report.csv},
  \texttt{ICR\_00\_csv\_conformite\_problems.csv} (optionnel),
  \texttt{ICR\_donnees\_preparees.csv}.

  \item \textbf{Comparaisons} (\texttt{comparisons/}) :
  \texttt{ICR\_resume\_tous\_sites.csv},
  \texttt{ICR\_comparaison\_completude\_sites.png},
  \texttt{ICR\_comparaison\_ir\_sites.png}.

  \item \textbf{Métriques} (\texttt{06\_metrics/}) :
  \texttt{ICR\_metrics\_pertinence\_tous\_sites.csv},
  \texttt{ICR\_metrics\_pertinence\_heatmap\_sites.png},
  \texttt{ICR\_metrics\_pertinence\_score\_sites.png}.
\end{itemize}

Logs : \texttt{ICR\_YYYYMMDD\_HHMM.log} (sous \texttt{PROJECT\_DIR/logs/}, hors de \texttt{results/ICR/}).

\subsubsection{Par site (\texttt{sites/<site\_name>/})}
Fichiers CSV par site : courbe temps-espèces, TEE/IR, fréquences,
résumé de site et métriques de pertinence, avec les graphiques PNG associés.

Sorties conditionnelles : \texttt{ICR\_05\_modeles.csv}, \texttt{ICR\_06\_occupation\_spatiale.csv}, \texttt{ICR\_05\_temporel\_vs\_spatial.png}, \texttt{ICR\_06\_CA\_placettes\_especes.png}.

\subsection{Flux d'exécution}
\begin{verbatim}
run_analysis(CONFIG)
  -> validate_config()
  -> resolve_input_file()
  -> read_delim_auto()
  -> audit_csv_conformity()
  -> export_csv_conformity_report()
  -> prepare_data()
  -> for each site: analyze_site()
       -> calculs cumul/TEE/frequences/spatial
       -> modeles lineaire/hyperbolique
       -> metriques scientifiques
       -> exports CSV/PNG
  -> consolidation multi-sites
  -> graphiques globaux
  -> log_output_manifest()
\end{verbatim}

\section{Robustesse et gestion d'erreurs}
\label{sec:robustesse}

\subsection{Niveaux de protection}
\begin{itemize}[leftmargin=1.2cm]
  \item \texttt{validate\_config()} : vérifie la présence et le type de chacune des clés de configuration (\texttt{stop()} explicite sinon), notamment la monotonie de \texttt{freq\_breaks}, la cohérence de longueur avec \texttt{freq\_labels} et \texttt{min\_visits\_for\_model >= 2}.
  \item \texttt{prepare\_data()} : contrôle du schéma et des valeurs manquantes/vides sur \texttt{site}/\texttt{visite\_id}/\texttt{espece}/\texttt{date} (\texttt{stop()} explicite, message ciblant la colonne en cause).
  \item \texttt{tryCatch} locaux dans \texttt{analyze\_site()} : chaque tentative d'ajustement (\texttt{fit\_linear\_model()}, \texttt{fit\_hyperbolic\_model()}) et d'analyse CA/AFC (\texttt{run\_ca()}) est isolée ; un échec est journalisé via \texttt{log\_warning()} et n'interrompt ni le site ni le pipeline.
  \item \texttt{tryCatch} sur chaque écriture CSV (\texttt{readr::write\_csv()}) : un échec d'écriture provoque un \texttt{stop()} explicite mentionnant le fichier concerné (contrairement aux échecs de modélisation, jugés non recouvrables).
  \item \textbf{Absence de \texttt{tryCatch} global} : à la différence du script CHEGD de référence (qui encapsule son point d'entrée principal), \texttt{run\_analysis()} n'est pas protégée globalement ; toute erreur de validation amont (configuration, CSV non conforme en mode strict, données invalides) interrompt l'exécution complète du script.
\end{itemize}

\subsection{Fallbacks notables}
\begin{itemize}[leftmargin=1.2cm]
  \item Package \texttt{minpack.lm} absent $\Rightarrow$ pas de modèle hyperbolique ni de complétude (valeurs \texttt{NA}), message \texttt{log\_info()} explicite au démarrage.
  \item Package \texttt{vegan} absent, ou colonne \texttt{placette} absente/vide $\Rightarrow$ pas de CA/AFC (\texttt{ICR\_06\_CA\_placettes\_especes.png} non généré).
  \item \texttt{nb\_visites < min\_visits\_for\_model} $\Rightarrow$ pas de modélisation pour le site concerné (\texttt{ICR\_05\_modeles.csv} non généré).
  \item Absence de colonne \texttt{placette} $\Rightarrow$ pas d'occupation spatiale (\texttt{occ\_tbl} \texttt{NULL}), pas de graphique spatio-temporel, et corrélation \texttt{spearman\_temporal\_spatial} non calculée.
\end{itemize}

\subsection{Politique de sévérité des anomalies}
\begin{longtable}{|p{2.0cm}|p{6.0cm}|p{6.3cm}|}
\hline
\textbf{Niveau} & \textbf{Déclencheurs typiques} & \textbf{Comportement effectif} \\
\hline
INFO & Package optionnel absent, mode d'exécution, volumétrie des données & Message informatif, aucune interruption \\
\hline
ALERTE & CSV non conforme en mode tolérant, échec d'ajustement de modèle, échec CA/AFC & \texttt{log\_warning()} (console + fichier + \texttt{stderr}), le pipeline se poursuit avec livrable(s) correspondant(s) marqué(s) \texttt{OPTIONNEL\_NON\_GENERE} \\
\hline
ERREUR & Colonne obligatoire absente, valeurs manquantes sur les champs clés, CSV non conforme en mode strict, échec d'écriture d'un CSV & \texttt{stop()} : arrêt immédiat de l'exécution \\
\hline
\end{longtable}

\section{Contrôles qualité et non-régression}

\subsection{Indicateurs d'assurance qualité (QA) réellement exportés}
\begin{itemize}[leftmargin=1.2cm]
  \item \texttt{ICR\_00\_csv\_conformite\_report.csv} : statut \texttt{CONFORME}/\texttt{NON\_CONFORME}, colonnes manquantes/supplémentaires, nombre de problèmes de parsing.
  \item \texttt{ICR\_00\_csv\_conformite\_problems.csv} (optionnel) : détail des problèmes \texttt{readr::problems()}.
  \item Manifeste des livrables (\texttt{log\_output\_manifest()}) : statuts \texttt{OK}/\texttt{MANQUANT}/\texttt{OPTIONNEL\_NON\_GENERE} par fichier attendu, journalisés en fin d'exécution.
  \item \texttt{class\_ir}/\texttt{class\_completude} dans \texttt{ICR\_07\_metrics\_pertinence.csv} : classes qualitatives directement exploitables pour un contrôle de cohérence.
\end{itemize}

\subsection{Critères techniques d'acceptation}
\begin{enumerate}
  \item Exécution sans erreur fatale pour un CSV conforme au schéma attendu.
  \item Génération complète des sorties CSV/figures attendues par site et au global (manifeste sans statut \texttt{MANQUANT} non justifié).
  \item Log horodaté présent sous \texttt{PROJECT\_DIR/logs/}.
  \item Cohérence des indicateurs TEE/Ir/complétude avec leurs seuils documentés (bornes $[0,1]$).
\end{enumerate}

\subsection{Stratégie de tests recommandée}
\begin{itemize}[leftmargin=1.2cm]
  \item \textbf{Tests unitaires} : \texttt{calc\_tee\_ir()}, \texttt{calc\_cumulative\_metrics()}, \texttt{calc\_species\_frequency()} et \texttt{calc\_scientific\_metrics()} sur des jeux synthétiques à richesse connue (valeurs de TEE/Ir/score calculables à la main).
  \item \textbf{Tests d'intégration} : exécution complète sur un jeu de données réduit, avec et sans colonne \texttt{placette}.
  \item \textbf{Tests de non-régression} : comparaison des CSV clés (\texttt{ICR\_04\_resume\_site.csv}, \texttt{ICR\_07\_metrics\_pertinence.csv}) avec des instantanés de référence versionnés.
  \item \textbf{Tests de robustesse} : dates ambiguës, doublons volontaires, sites avec un nombre de visites proche ou inférieur à \texttt{min\_visits\_for\_model}, absence simulée de \texttt{minpack.lm}/\texttt{vegan}.
\end{itemize}

\subsection{Matrice de validation des artefacts}
\begin{longtable}{|p{3.5cm}|p{5.2cm}|p{5.6cm}|}
\hline
\textbf{Artefact} & \textbf{Règle de validation} & \textbf{Preuve attendue} \\
\hline
Ingestion (\texttt{ICR\_donnees\_preparees.csv}) & Aucune valeur manquante sur \texttt{site}/\texttt{visite\_id}/\texttt{espece}/\texttt{date} & Comptage avant/après déduplication journalisé \\
\hline
Courbe temps-espèces (\texttt{ICR\_01}) & \texttt{nb\_cumule} croissant (ou stable) avec \texttt{visite\_index} & Contrôle de monotonie \\
\hline
TEE/Ir (\texttt{ICR\_02}) & $TEE_i, Ir_i \in [0,1]$ pour toute visite $i$ & Bornage vérifiable directement sur le CSV \\
\hline
Fréquences (\texttt{ICR\_03}) & Classes cohérentes avec \texttt{freq\_breaks}/\texttt{freq\_labels} & Recoupement \texttt{cut()} manuel sur \texttt{prop\_visites} \\
\hline
Modèles (\texttt{ICR\_05}, si activés) & $R^2 \in [0,1]$, asymptote $> \max(\texttt{nb\_cumule})$ & Comparaison directe aux données observées \\
\hline
Pertinence scientifique (\texttt{ICR\_07}) & $\texttt{score\_pertinence} \in [0,1]$, classes \texttt{class\_ir}/\texttt{class\_completude} cohérentes avec les seuils documentés & Recalcul manuel de la moyenne pondérée \\
\hline
\end{longtable}

\subsection{Menaces à la validité}
\begin{itemize}[leftmargin=1.2cm]
  \item \textbf{Interne} : dépendance forte à la qualité de saisie de \texttt{visite\_id}/\texttt{date} (une erreur de saisie fusionne ou scinde artificiellement des visites).
  \item \textbf{Externe} : transférabilité limitée si le protocole d'échantillonnage (nombre de visites, granularité de \texttt{placette}) diffère fortement d'un site à l'autre ; les comparaisons inter-sites doivent en tenir compte.
  \item \textbf{De conclusion} : prudence sur la complétude/Ir des sites dont le nombre de visites est proche du seuil \texttt{min\_visits\_for\_model}, où les estimations restent statistiquement fragiles.
\end{itemize}

\section{Complexité et performances (qualitatives)}
\begin{itemize}[leftmargin=1.2cm]
  \item \texttt{calc\_cumulative\_metrics()}/\texttt{calc\_tee\_ir()} : implémentées via une boucle explicite sur les visites (\texttt{for} sur \texttt{seq\_len(nrow(visits\_tbl))}) avec un filtrage cumulatif à chaque itération, soit un coût approximativement $O(v^2)$ par site dans le pire cas ($v$ = nombre de visites) — acceptable pour les volumes usuels d'un inventaire de terrain, mais à surveiller pour des séries de visites très longues.
  \item Le commentaire d'en-tête du script (Phase 2, v1.2) mentionne une \og{}vectorisation des calculs cumulés\fg{} ; l'implémentation actuelle conserve toutefois une boucle explicite par visite pour ces deux fonctions — écart entre intention documentée et code effectif à conserver à l'esprit lors d'une future refactorisation.
  \item \texttt{calc\_species\_frequency()}/\texttt{calc\_spatial\_occupancy()} : agrégations vectorisées \texttt{dplyr}, coût proche de $O(n)$ sur le nombre d'observations.
  \item \texttt{fit\_hyperbolic\_model()} : coût dominé par les itérations de Levenberg-Marquardt (500 au maximum), généralement négligeable devant les boucles ci-dessus pour les tailles de site habituelles.
  \item \texttt{BENCHMARK\_MODE} : \texttt{bench\_time()} chronomètre \texttt{calc\_cumulative\_metrics()}, \texttt{calc\_tee\_ir()} et \texttt{calc\_species\_frequency()} par site lorsqu'activé.
\end{itemize}

\subsection{Pistes d'optimisation}
\begin{itemize}[leftmargin=1.2cm]
  \item Vectoriser \texttt{calc\_cumulative\_metrics()}/\texttt{calc\_tee\_ir()} (cumul par opérations ensemblistes vectorisées plutôt que boucle explicite) pour les sites à très grand nombre de visites.
  \item Réutiliser \texttt{freq\_tbl}/\texttt{occ\_tbl} déjà calculées comme entrées directes de \texttt{calc\_scientific\_metrics()} sans recalcul intermédiaire.
  \item Limiter la recréation des répertoires thématiques réservés lorsqu'ils restent vides d'une exécution à l'autre.
\end{itemize}

\section{Sécurité, conformité et exploitation}

\subsection{Sécurité des données}
Le script opère exclusivement sur des fichiers locaux (aucun appel réseau, aucune
transmission de données). Les journaux et les CSV générés peuvent contenir des noms de
sites potentiellement sensibles (localisation d'espèces protégées) : leur diffusion doit
suivre la politique de confidentialité du projet.

\subsection{Conformité et bonnes pratiques}
\begin{itemize}[leftmargin=1.2cm]
  \item Aucune dépendance à des secrets ou identifiants (pas de connexion externe).
  \item Chemins résolus de manière relative au projet (\texttt{resolve\_path()}), limitant les risques de chemins absolus non maîtrisés.
  \item Traçabilité des sorties via le manifeste (\texttt{log\_output\_manifest()}) et les logs horodatés.
\end{itemize}

\subsection{Procédure opératoire standard (POS)}
\begin{enumerate}
  \item Vérifier la présence et le schéma du fichier d'entrée (\texttt{data/observations.csv}).
  \item Exécuter le script (\texttt{Rscript} ou session interactive).
  \item Contrôler le log généré et le rapport de conformité CSV.
  \item Diffuser les exports CSV/PNG validés aux parties prenantes.
  \item Archiver le run (log + rapport de conformité + CSV consolidés) pour traçabilité.
\end{enumerate}

\section{Réplicabilité et transparence scientifique}
Le script ne génère pas actuellement de manifeste de run structuré (type JSON) ; seul le
manifeste texte des livrables (\texttt{log\_output\_manifest()}) et le fichier log horodaté
assurent la traçabilité. Pour renforcer la réplicabilité inter-campagnes, un manifeste de
run \texttt{run\_manifest.json} pourrait consigner :
\begin{itemize}[leftmargin=1.2cm]
  \item \texttt{run\_id}, \texttt{timestamp\_start}, \texttt{timestamp\_end} ;
  \item \texttt{r\_version}, versions des packages \texttt{dplyr}/\texttt{ggplot2}/\texttt{minpack.lm}/\texttt{vegan} ;
  \item \texttt{input\_file}, empreinte (hash) du fichier d'entrée ;
  \item \texttt{config\_effective} (valeurs réelles de \texttt{CONFIG} après surcharge par variables d'environnement) ;
  \item \texttt{outputs\_generated} (liste des fichiers effectivement produits) ;
  \item \texttt{qa\_summary} (statut de conformité CSV, nombre de sites traités).
\end{itemize}

\section{Discussion}

\subsection{Améliorations techniques (v1.4)}
\subsubsection{Phase 1 : Fiabilisation}
\begin{itemize}
  \item D\'etection automatique du d\'elimiteur (CSV/TSV/TXT)
  \item Parsing flexible des dates via format configurable
  \item D\'eduplication multi-cl\'e : \texttt{(site, placette, date, visite\_id, espece)}
  \item Gestion granulaire des erreurs CSV (colonne manquante, type invalide, etc.)
\end{itemize}

\subsubsection{Phase 2 : Performance}
\begin{itemize}
  \item Optimisation du modèle hyperbolique : \texttt{minpack.lm::nlsLM} avec détection de convergence
  \item Calculs cumulés intentionnellement documentés comme \og{}vectorisés\fg{} dans l'en-tête du script ; l'implémentation effective conserve une boucle explicite par visite (voir \S\,\ref{sec:robustesse} et la section Complexité et performances)
  \item M\'emorisation des graphiques ggplot2 avant export PNG
  \item Cache des analyses de CA/AFC conditionnelle
\end{itemize}

\subsubsection{Phase 3 : Graphiques et visualisations}
\begin{itemize}
  \item Cartes de chaleur multi-sites (heatmap) avec échelle de couleur
  \item Tableaux de bord par site (7 graphiques dans une grille)
  \item Comparaisons visuelles de complétude (barplot multi-sites)
  \item Analyses de CA/AFC si \texttt{vegan} disponible et placette présente
\end{itemize}

\subsubsection{Phase 4 : Industrialisation}
\begin{enumerate}
  \item \textbf{Configuration centralisée} : \texttt{get\_embedded\_config()} avec 14+ paramètres
  \item \textbf{Logging systématisé} : \texttt{setup\_logging()} + fonctions log\_*() avec fichiers horodatés
  \item \textbf{Organisation thématique} : hiérarchie logique /sites/, /06\_metrics/, /00\_data\_quality/, /comparisons/
  \item \textbf{Métriques avancées} : calcul de pertinence scientifique, dashboards analytiques, manifeste final
\end{enumerate}

\subsection{Limites connues}
\begin{itemize}
  \item Sans \texttt{placette}, pas d'analyse spatiale ni de CA/AFC, et la corrélation temporel/spatial de \texttt{calc\_scientific\_metrics()} n'est pas calculée.
  \item Sans \texttt{minpack.lm}, pas d'asymptote hyperbolique ni de complétude associée ; le score de pertinence se rabat alors sur le $R^2$ du modèle linéaire.
  \item Si \texttt{nb\_visites < min\_visits\_for\_model}, pas de modélisation pour le site (impact direct sur \texttt{completude} et \texttt{score\_pertinence}).
  \item En mode strict (défaut), un CSV non conforme bloque tout le pipeline sans possibilité de reprise partielle.
  \item Absence de \texttt{tryCatch} global autour de \texttt{run\_analysis()} : une erreur de configuration ou de données en amont interrompt l'exécution complète plutôt que de permettre un traitement partiel des sites valides.
  \item Les clés \texttt{freq\_breaks}, \texttt{freq\_labels}, \texttt{make\_ca}, \texttt{min\_visits\_for\_model}, \texttt{width}, \texttt{height} et \texttt{dpi} ne sont pas surchargeables par variable d'environnement, contrairement à ce que suggère le principe de configuration centralisée.
\end{itemize}

\section{Protocole de référence}
Les indicateurs implémentés (courbe temps-espèces, taux d'espèces exclusives, indice de
représentativité) s'appuient sur les principes présentés dans le cours \og{}Cours 24
Inventaires\fg{} du Diplôme Universitaire de Mycologie 2026 (Professeur Pierre-Arthur
Moreau, Université de Lille, UFR3S PHAR, diapositives 25 à 48), qui formalise la logique
d'accumulation d'espèces au fil des visites et l'estimation de la représentativité
d'un inventaire de terrain. Contrairement au protocole CHEGD de Sellier (module fiabilité
du script de référence), aucun protocole normatif externe spécifique n'est requis ici : le
script est entièrement autonome et ne dépend d'aucun classeur ou référentiel externe.

\section{Évolutions techniques recommandées}
\begin{itemize}[leftmargin=1.2cm]
  \item Ajouter un \texttt{tryCatch} global autour de \texttt{run\_analysis()} pour aligner la robustesse sur le script CHEGD de référence et éviter l'arrêt complet sur une erreur isolée.
  \item Vectoriser \texttt{calc\_cumulative\_metrics()}/\texttt{calc\_tee\_ir()} pour supprimer la boucle explicite par visite et réduire la complexité quadratique observée.
  \item Exposer \texttt{freq\_breaks}/\texttt{freq\_labels}/\texttt{min\_visits\_for\_model} par variable d'environnement, en cohérence avec le principe de configuration centralisée déjà appliqué aux autres clés.
  \item Générer un manifeste de run structuré (\texttt{run\_manifest.json}) pour la traçabilité inter-campagnes (voir \S\,Réplicabilité et transparence scientifique).
  \item Ajouter une suite de tests unitaires sur les formules métier (TEE, Ir, complétude, score de pertinence) afin de sécuriser les futures refactorisations.
\end{itemize}

\section{Références bibliographiques}
\begin{itemize}[leftmargin=1.2cm]
  \item Moreau, P.-A. \textit{Cours 24 --- Inventaires}. Diplôme Universitaire de Mycologie 2026, Université de Lille (UFR3S PHAR), diapositives 25--48.
  \item Colwell, R.K. \& Coddington, J.A. (1994). Estimating terrestrial biodiversity through extrapolation. \textit{Philosophical Transactions of the Royal Society B}, 345(1311), 101--118.
  \item Wickham, H. (2016). \textit{ggplot2: Elegant Graphics for Data Analysis}. Springer.
  \item Wickham, H., \& Grolemund, G. (2017). \textit{R for Data Science}. O'Reilly.
  \item Oksanen, J. et al. \textit{vegan: Community Ecology Package}. Paquet R (analyse de correspondance CA/AFC).
  \item Elzhov, T.V., Mullen, K.M., Spiess, A.-N., \& Bolker, B. \textit{minpack.lm: R Interface to the Levenberg-Marquardt Nonlinear Least-Squares Algorithm}. Paquet R.
\end{itemize}

% ============================================================================================================================
\section*{Contact et versionnement}
\addcontentsline{toc}{section}{Contact et versionnement}
% ============================================================================================================================

\begin{table}[htbp]
  \centering
  \begin{tabular}{ll}
    \toprule
    \textbf{Auteur}          & Eddy Boite \\
    \textbf{Dépôt GitHub}    & \texttt{eddyboite-nat/statistics} \\
    \textbf{Branche}         & \texttt{main} \\
    \textbf{Version script}  & 1.4 \\
    \textbf{Date de ce document} & 9 juillet 2026 \\
    \bottomrule
  \end{tabular}
\end{table}

Pour toute évolution du script, ouvrir une issue dans le dépôt GitHub
ou documenter les changements dans le journal \texttt{CHANGELOG.md} du projet.

% ============================================================================================================================

\end{document}
